\chapter{Linear and Polynomial Regression}
\labinfo{02\_Linear\_Regression}{linear\_regression.py}{UCI Red Wine Quality: 1,599 wines; alcohol, volatile acidity, sulphates, and density}{Regression}{Degree-2 test RMSE 0.632; $R^2=0.325$}

\section{Method}
Wine quality (scores 3--8) was predicted from four selected measurements. The
script compares simple and multiple least-squares regression with polynomial
models of degrees 1--6. Degree was selected by five-fold CV on the training
split; input standardization is learned from training data only.

\section{Results}
Degree 2 gave the lowest CV RMSE (0.661). Its held-out RMSE was 0.632 and
$R^2=0.325$, slightly better than the multiple linear model (RMSE 0.642,
$R^2=0.304$). Higher degrees overfit: degree 6 had CV RMSE 20.941. The modest
$R^2$ indicates that these four measurements explain only part of the rating.

\begin{figure}[H]\centering
\includegraphics[width=.82\linewidth]{\figpath{02_Linear_Regression}{degree_vs_rmse.pdf}}
\caption{Training, cross-validation, and test errors by polynomial degree.}
\end{figure}

\section{Python source}
\lstinputlisting[caption={Complete regression implementation},label={lst:lab02}]{02_Linear_Regression/linear_regression.py}
